\Article{64}
Gehen wir von der letzten im vorigen Artikel gefundenen Regel aus, so finden wir
\[
\begin{array}{c|l}
\text{Zahl}&\text{Charakter}\equiv\\\hline
-1+i&\frac18(aa+2ab-bb-1)\\
-1-i&\frac18(-aa+2ab+bb+1)\\
+1-i&\frac18(aa+2ab+3bb-1)
\end{array}
\]

Das folgt sogleich daraus, dass der Charakter von $i$ gleich $\frac14(aa+bb-1)$, der von $-1$ dagegen $\frac12(aa+bb-1)\equiv\frac12bb$ ist, da $aa-1$ stets die Form $8n$ hat. Diese vier Regeln sind offenbar, obwohl sie bisher aus der Untersuchung von Zahlenbeispielen gewonnen wurden, so miteinander verbunden, dass mit dem abgeschlossenen Beweis einer von ihnen zugleich die drei übrigen bewiesen sind. Es ist kaum nötig zu bemerken, dass auch in diesen Regeln nur $a$ als ungerade und $b$ als gerade vorausgesetzt wird.

Will man Formeln verwenden, die auf primäre Moduln beschränkt sind, so kann man die folgende Form benutzen:
\[
\begin{array}{c|l}
\text{Zahl}&\text{Charakter}\equiv\\\hline
-1+i&\frac14(-a-b+1-bb)\\
-1-i&\frac14(a-b-1+bb)\\
+1-i&\frac14(-a-b+1+bb)
\end{array}
\]

Die einfachsten Formeln ergeben sich, wenn wir wie zu Beginn unserer Untersuchung von Zahlenbeispielen Moduln der ersten und zweiten Art unterscheiden. Die Charaktere sind
\[
\begin{array}{c|c|c}
\text{Zahl}&\text{Für Moduln der ersten Art}&\text{Für Moduln der zweiten Art}\\\hline
-1+i&\frac14(-a-b+1)&\frac14(-a-b-3)\\
-1-i&\frac14(a-b-1)&\frac14(a-b+3)\\
+1-i&\frac14(-a-b+1)&\frac14(-a-b+5)
\end{array}
\]

\Needspace{14\baselineskip}
\Article{65}
Für die Zahl $-1+2i$, zu der wir nun übergehen, verwenden wir ebenfalls dieselbe Unterscheidung zwischen den Moduln $a+bi$, für die $a\equiv1,b\equiv0$ gilt, und denen, für die $a\equiv3,b\equiv2$ gilt. Die Tabelle in Artikel~62 zeigt für jene Zahl die folgenden Zuordnungen:
\SourceLeaf{147}%
\[
\begin{array}{c|l}
\text{Charakter}&\text{Moduln der ersten Art}\\\hline
0&-3+8i,\ +5-8i,\ +9+4i,\ -11+4i\\
1&+1+4i,\ +5-4i,\ -7,\ -3-8i\\
2&+1-4i,\ +5+8i,\ -7-8i,\ -11\\
3&-3,\ +5+4i,\ +9-4i,\ -7+8i,\ -11-4i
\end{array}
\]

Führen wir jeden dieser Moduln auf seinen absolut kleinsten Rest nach dem Modul $-1+2i$ zurück, so sehen wir, dass alle, denen der Charakter $0$ entspricht, $\equiv1$ sind; diejenigen, denen der Charakter $1$ entspricht, sind $\equiv i$; diejenigen, deren Charakter $2$ ist, werden $\equiv-1$; schließlich werden alle, deren Charakter $3$ ist, $\equiv-i$. Die Charaktere der Zahlen $1,i,-1,-i$ für den Modul $-1+2i$ sind aber selbst der Reihe nach $0,1,2,3$. Somit ist in allen diesen $17$ Beispielen der Charakter der Zahl $-1+2i$ bezüglich eines Moduls $a+bi$ der ersten Art mit dem Charakter dieser Zahl bezüglich des Moduls $-1+2i$ identisch.

Ebenso findet man mithilfe der Tabelle die Zuordnungen
\[
\begin{array}{c|l}
\text{Charakter}&\text{Moduln der zweiten Art}\\\hline
0&+3+2i,\ -5-2i,\ -1+10i,\ -1-10i,\ +11+6i\\
1&+3-2i,\ -1+6i,\ -5-6i,\ +7+10i,\ +7-10i\\
2&-5+2i,\ -1-6i,\ +7-2i\\
3&-1-2i,\ +7+2i,\ -5+6i,\ +3+10i,\ +3-10i,\ +11-6i
\end{array}
\]

Führt man diese Moduln auf ihre kleinsten Reste nach dem Modul $-1+2i$ zurück, so erweisen sich alle, denen der Reihe nach die Charaktere $0,1,2,3$ entsprechen, den Zahlen $-1,-i,+1,+i$ kongruent. Diese Zahlen selbst haben aber, wenn umgekehrt $-1+2i$ als Modul genommen wird, der Reihe nach die Charaktere $2,3,0,1$. Somit unterscheidet sich in allen diesen $19$ Beispielen der Charakter der Zahl $-1+2i$ bezüglich eines Moduls der zweiten Art um zwei Einheiten vom Charakter dieser Zahl bezüglich der als Modul genommenen Zahl $-1+2i$.

Ferner sieht man ohne Mühe, dass für die Zahl $-1-2i$ ganz Ähnliches gilt.

\Article{66}
Für die Zahl $-3$ lassen wir die Unterscheidung zwischen Moduln der ersten und zweiten Art weg, da das Ergebnis zeigt, dass sie hier überflüssig ist. Es gelten somit die Zuordnungen
\SourceLeaf{148}%
\[
\begin{array}{c|l}
\text{Charakter}&\text{Moduln}\\\hline
0&\begin{gathered}-1+6i,\ -1-6i,\ -7,\ -5+6i,\\-5-6i,\ -11,\ 11+6i,\ 11-6i\end{gathered}\\[3pt]
1&\begin{gathered}-1-2i,\ 1-4i,\ -5+2i,\ 5+4i,\ 7+2i,\\5-8i,\ -1+10i,\ -7-8i,\ -11-4i,\ 7-10i\end{gathered}\\[3pt]
2&\begin{gathered}3+2i,\ 3-2i,\ -3+8i,\ -3-8i,\\9+4i,\ 3+10i,\ 3-10i\end{gathered}\\[3pt]
3&\begin{gathered}-1+2i,\ 1+4i,\ -5-2i,\ 5-4i,\ 7-2i,\\5+8i,\ -1-10i,\ -7+8i,\ -11+4i,\ 7+10i\end{gathered}
\end{array}
\]

\Needspace{12\baselineskip}
Führen wir diese Moduln auf ihre kleinsten Reste nach dem Modul $3$ zurück, so sehen wir, dass diejenigen, denen der Charakter $0$ entspricht, teils $\equiv1$, teils $\equiv-1$ sind; diejenigen, deren Charakter $1$ ist, werden entweder $\equiv1-i$ oder $\equiv-1+i$; diejenigen, deren Charakter $2$ ist, werden entweder $\equiv i$ oder $\equiv-i$; schließlich sind diejenigen, denen der Charakter $3$ zukommt, entweder $\equiv1+i$ oder $\equiv-1-i$. Aus dieser Untersuchung von Zahlenbeispielen schließen wir daher, dass der Charakter der Zahl $-3$ für einen Modul, der eine unter ihren Assoziierten primäre Primzahl ist, mit dem Charakter dieser Zahl selbst identisch ist, wenn $3$ oder gleichbedeutend $-3$ als Modul betrachtet wird.

\newpage
\Article{67}
Eine ähnliche Untersuchung von Zahlenbeispielen für andere Primzahlen zeigt, dass die Zahlen $3\pm2i,-1\pm6i,7\pm2i,-5\pm6i$ usw. Sätze liefern, die dem in Artikel~65 für die Zahl $-1+2i$ gefundenen ähnlich sind, dass sich dagegen die Zahlen $1\pm4i,5\pm4i,-3\pm8i,5\pm8i,9\pm4i$ usw. ebenso verhalten wie die Zahl $-3$. Die Untersuchung von Zahlenbeispielen führt somit zu einem höchst eleganten Satz,. den man entsprechend der Theorie quadratischer Reste in der Arithmetik reeller Zahlen den \textsc{Fundamentalsatz} der Theorie biquadratischer Reste nennen darf, nämlich:

\emph{Bezeichnen $a+bi,a'+b'i$ verschiedene, unter ihren jeweiligen Assoziierten primäre Primzahlen, also solche, die der Einheit nach dem Modul $2+2i$ kongruent sind, so ist der biquadratische Charakter der Zahl $a+bi$ bezüglich des Moduls $a'+b'i$ mit dem Charakter der Zahl $a'+b'i$ bezüglich des Moduls $a+bi$ identisch, wenn entweder beide Zahlen $a+bi,a'+b'i$ oder wenigstens eine von ihnen zur ersten Art gehören, also nach dem Modul $4$ der Einheit kongruent sind. Dagegen unterscheiden sich jene Charaktere um zwei Einheiten, wenn keine der Zahlen $a+bi,a'+b'i$ zur ersten Art gehört, also wenn beide nach dem Modul $4$ der Zahl $3+2i$ kongruent sind.}

\SourceLeaf{149}%
Trotz der großen Einfachheit dieses Satzes zählt sein Beweis jedoch zu den tiefsten Geheimnissen der höheren Arithmetik. Zumindest beim gegenwärtigen Stand lässt er sich nur durch die subtilsten Untersuchungen erschließen, die weit über die Grenzen der vorliegenden Abhandlung hinausgingen. Daher behalten wir uns die Veröffentlichung dieses Beweises sowie die Entwicklung des Zusammenhangs zwischen diesem Satz und jenen, die wir am Anfang dieser Abhandlung durch die Untersuchung von Zahlenbeispielen festzustellen begonnen haben, für eine dritte Abhandlung vor. Zum Abschluss geben wir jedoch schon hier an, was zum Beweis der in Artikel~63.~64 vorgeschlagenen Sätze erforderlich ist.

\Needspace{15\baselineskip}
\Article{68}
Wir beginnen mit Primzahlen $a+bi$, für die $b=0$ gilt (der dritten Art in Artikel~34). Damit die Zahl unter ihren Assoziierten primär ist, muss hier also $a$ eine negative reelle Primzahl der Form $-(4n+3)$ sein, für die wir $-q$ schreiben; solche Zahlen sind $-3,-7,-11,-19$ usw. Bezeichnet $\lambda$ den Charakter der Zahl $1+i$, wenn jene Zahl als Modul genommen wird, so muss gelten
\[
i^\lambda\equiv(1+i)^{\frac14(qq-1)}\equiv2^{\frac18(qq-1)}\pdot i^{\frac18(qq-1)}\pmod q
\]
Bekanntlich ist aber $2$ ein quadratischer Rest oder Nichtrest von $q$, je nachdem, ob $q$ die Form $8n+7$ oder die Form $8n+3$ hat. Daraus schließen wir, dass allgemein
\[
2^{\frac12(q-1)}\equiv(-1)^{\frac14(q+1)}\equiv i^{\frac12(q+1)}\pmod q
\]
gilt, und somit folgt durch Erheben zur Potenz mit dem Exponenten $\frac14(q+1)$
\[
2^{\frac18(qq-1)}\equiv i^{\frac18(q+1)^2}\pmod q
\]
Die vorhergehende Gleichung nimmt daher die Form
\[
i^\lambda\equiv i^{\frac18(q+1)^2+\frac18(qq-1)}\equiv i^{\frac14(qq+q)}\pmod q
\]
an, woraus
\[
\lambda\equiv\frac14(qq+q)\equiv\frac14(q+1)^2-\frac14(q+1)\pmod4
\]
folgt; oder, da $\frac14(q+1)^2\equiv0\pmod4$ gilt, $\lambda\equiv-\frac14(q+1)\equiv\frac14(a-1)\pmod4$. Das ist gerade der Satz von Artikel~63 für den Fall $b=0$.

\SourceLeaf{150}%
